% Guia do professor — Capítulo 11: Circulação Geral
\section{Capítulo 11 --- Circulação Geral}

\subsection*{Visão geral e ritmo}

O capítulo que conecta o micro (Caps.~1--10) ao planeta: desequilíbrio
radiativo, células, jatos (momento angular e vento térmico) e ondas de
Rossby. O capítulo do livro é longo (60\,pp); estas notas destilam a
espinha dorsal física. \textbf{Ritmo: 2 aulas de 2\,h.}

\begin{itemize}
  \item \textbf{Aula 1} — os dois perfis radiativos e o transporte exigido
        (integral no quadro), Hadley, momento angular e o jato subtropical,
        o mapa-múndi climático. Casos~1--2 (15\,min; o EBM rende).
  \item \textbf{Aula 2} — vento térmico (dedução completa), jato polar,
        Rossby ($\beta$ como restauração), cinturões e bloqueios.
        Casos~3--4 (15\,min).
\end{itemize}

\subsection*{Objetivos de aprendizagem}

\begin{enumerate}
  \item calcular o transporte meridional exigido a partir dos perfis
        radiativos;
  \item explicar Hadley, seus limites ($\sim$30°) e o jato subtropical por
        momento angular;
  \item deduzir e aplicar o vento térmico; localizar o jato polar;
  \item explicar a restauração $\beta$ e usar $c = U - \beta/k^2$;
  \item ler o mapa-múndi climático como consequência das células.
\end{enumerate}

\subsection*{Pré-requisitos a reativar}

Insolação por latitude (Cap.~2, Caso~2!); geostrofia e hidrostática
(Caps.~1, 10); momento angular (mecânica); vorticidade (definir
rapidamente — físicos aceitam $\zeta = \nabla\times\vec U$ de imediato).

\subsection*{Armadilhas conceituais frequentes}

\begin{itemize}
  \item \textbf{``Hadley vai até o polo''}: o modelo de uma célula falha
        por momento angular (134\,m/s!) e baroclinia; T2 quantifica.
  \item \textbf{Jato como ``rio de ar'' isolado}: é a expressão vertical do
        contraste térmico (vento térmico), não uma entidade à parte.
  \item \textbf{Rossby para leste}: a FASE das ondas longas anda para
        oeste em relação ao escoamento; a advecção pelo $U$ costuma vencer
        — distinguir fase × grupo × advecção.
  \item \textbf{ZCIT no equador geográfico}: migra com a estação e fica em
        média ao norte — vale mostrar a climatologia.
  \item \textbf{Sinais no HS}: vento térmico e Buys-Ballot de novo com
        $f<0$; insistir nos desenhos com o sul para baixo.
\end{itemize}

\subsection*{Demonstração de baixo custo}

Tanque girante improvisado (balde em prato giratório com água + corante +
gelo no centro): células e ondas aparecem. Alternativa: vídeos do
laboratório de tanque do MIT (Marshall) — Hadley e baroclinia em 3\,min.

\subsection*{Problemas sugeridos do livro (exercícios do Cap.~11)}

Aquecimento: $\beta$ e período de Rossby em várias latitudes; vento térmico
com gradientes dados. Aprofundamento: transporte exigido de tabelas
radiativas; jato por camadas. Síntese: ``o que acontece com os cinturões
num planeta girando 2$\times$ mais rápido?''.
\emph{Confira a numeração na sua cópia (v1.02b).}

\subsection*{Uso do notebook \texttt{cap11\_circulacao.ipynb}}

\begin{center}
\begin{tabular}{lll}
\hline
Caso & Conteúdo & Momento sugerido \\
\hline
1 & Transporte exigido (os 5\,PW) & Aula 1, após a integral \\
2 & EBM meridional (\texttt{climlab}): $T(\phi)$, linha do gelo & Aula 1 \\
3 & Dispersão de Rossby e onda estacionária & Aula 2, após Rossby \\
4 & Vento térmico $\to$ jato num perfil realista & Aula 2, após dedução \\
\hline
\end{tabular}
\end{center}

O Caso~2 reaproveita o \texttt{climlab} do Cap.~2 em modo meridional — e
antecipa o gelo-albedo do Cap.~21.

\subsection*{Exercícios complementares e soluções}

Nas notas: T1--T5 e N1--N4. Soluções em \texttt{cap11\_solucoes.ipynb}
(\emph{não distribuir}): T3 com \texttt{sympy}; N2 é a ponte direta para o
Cap.~21. Pares: T1$\leftrightarrow$N1, T5$\leftrightarrow$N3.

\subsection*{Conexões com o resto do curso}

Insolação (Cap.~2); geostrofia/vorticidade (Cap.~10); ciclones nascem nas
ondas (Cap.~13); jato e CAT (Cap.~5); bloqueios e extremos (Cap.~21);
correntes oceânicas e ENSO (Cap.~21).
